% =============================================================================
% Chapter 3  The Physics-Informed Effective Hamiltonian Method  --  ~4 pp
% rubric: Introduction & Literature Review /20 (body-located: this chapter
% discharges the HD gate, "critical assessment").
% Plan: ThesisPlanning.md §5.3 (decision T-13: prior art, judged).
%
% PURPOSE     Present Wu et al.'s method at journal-review depth, in the register
%             of ATTRIBUTED PRIOR ART, then judge it. The judgement poses the three
%             questions Chapters 5, 6 and 7 answer.
% CRITERION   The opening sentence says this chapter presents reported prior art;
%             the closing section reads as a gap statement.
% SOURCES     Planning.md §3.1 (model, eta), §3.2 (operator dictionary), §3.3 (the
%             standard form as published), §3.6 (doubled space), §3.7 ("Standard
%             form, as published"), §3.8 (the paper's QSVT and time rule), §7.6
%             (the paper's Eqs 36-39); claims A-01 to A-13 stated neutrally here
%             and tested in §8.2 (sec:results-reproduction).
% WRITE WHEN  Now (wave W1; gate G1 passed).
%
% OUR measurements of the paper (the reproduction, the rebuilt circuits, the
% eigh/SVD finding, the measured conditioning) belong in Chapters 5 and 8, never
% in §3.2-3.5 (T-13). The paragraphs of the pre-plan draft that reported them,
% the legacy ladder reconstruction and the old "encoding wall" assessment are
% in the snapshot commit 2cd640c (2_body/3_pihm_collocation.tex).
% =============================================================================

\subsection{Overview}
\label{sec:pihm-overview}
% ~0.3 pp
% LAND, in order:
%   1. One opening sentence: this chapter reports Wu et al.'s method as published,
%      in the register of reviewed literature; this thesis's reproduction of it
%      is in sec:results-reproduction.
%   2. The three-ingredient move (the paragraph below, kept from the draft).
%   3. What the rest of the chapter does, and why it matters to Chapters 5-7.

Wu et al.\ recast solving a differential equation as three linked moves \cite{wu2025pihm}: represent a candidate solution as a quantum state carrying its Chebyshev coefficients; write the differential equation and every boundary or initial condition as linear operators that must annihilate that state; and assemble a Hamiltonian from those operators whose ground state is exactly the state satisfying all of them at once.

\subsection{The Latent Chebyshev Model}
\label{sec:pihm-model}
% ~0.5 pp
% LAND: the normalised basis; f_q(x) = <tau(x)|psi>; the scale eta; raw tau in
%   the rank-one rows. The closed form eta_e = ||b||^2 = (N/pi) int f^2/sqrt(1-x^2)
%   is OURS, so it is stated in §8.2 (eq:eta-e), not here.

\paragraph{The Chebyshev feature map.} A solution $f$ is approximated by a degree-$(N-1)$ truncated Chebyshev series, $N = 2^n$, carried by an $n$-qubit state
\begin{equation}
    \ket{\tau(x)}_n = \frac{1}{2^{n/2}}T_0(x)\ket{0}^{\otimes n} + \frac{1}{2^{(n-1)/2}}\sum_{k=1}^{N-1} T_k(x)\ket{k},
    \label{eq:pihm-tau}
\end{equation}
so that a coefficient vector $\ket{\psi}_n \in \mathbb{C}^N$ (unit norm, $\braket{\psi|\psi}=1$) represents the model
\begin{equation}
    f_Q(x) = \sqrt{\eta}\,\braket{\tau(x)|\psi}_n
    \label{eq:pihm-fq}
\end{equation}
for some positive scale $\eta$. Separating the \emph{direction} $\ket{\psi}$ from the \emph{scale} $\sqrt{\eta}$ this way matters because everything from here to the ground-state eigenvalue problem below depends only on $\ket{\psi}$: $\eta$ cancels out of every constraint $A\ket{\psi}=0$ and out of the Hamiltonian built from such constraints, so it plays no role until \Cref{eq:pihm-regular} fixes it, once, from a single known data point, entirely after the ground state has been found. Every equation from here on is written for $\ket{\psi}$ alone, with $\eta$ suppressed accordingly.

\subsection{The Residual, the Hamiltonian and Its Constraints}
\label{sec:pihm-residual}
% ~1.2 pp
% LAND, in order:
%   1. The derivative operator and its density (kept below).
%   2. The residual A_std = sum_j M_{a_j} G^j - M_r D0(x_s) and H_std (kept, plus
%      the source row, eq:pihm-A-source).
%   3. Least squares, with an exact kernel only for polynomial solutions.
%   4. Invariant versus regular constraints; k - 1 invariants for order k.
%   5. Variable coefficients through the lifts M_{x^p}; sources through
%      Maclaurin truncation (the paper's choice; its effect is measured in §8.2).
%   6. PDE invariants hand-picked from the known solution (heat: f_x(t, +-3/4) = 0;
%      wave: f(1/8, x) = f(7/8, x) = 0), the initial condition itself not a
%      constraint.
% OFFLOAD: the worked N = 4 instance and the lift entries -> Appendix C
%   (app:pihm-residual): "the worked N = 4 residual for the paper's {1,5,400}
%   pilot equation, and the entries of the lifts, are given in Appendix C".
% Convention (T-16): \G acts on coefficients, f -> f'. Wu et al. print this
%   matrix as G_n^T; say so once, here.

\paragraph{The derivative operator.} Differentiating a Chebyshev series produces another Chebyshev series whose coefficients are a fixed linear function of the original ones \cite{trefethen2000spectral}: writing this map in the normalised basis of \Cref{eq:pihm-tau} gives an $N \times N$ matrix $\G$ with
\begin{equation}
    \G_{k,p} =
    \begin{cases}
        2p, & p > k \ge 1,\ p-k\ \text{odd}, \\
        \sqrt{2}\,p, & k = 0,\ p\ \text{odd},\ p \ge 1, \\
        0, & \text{otherwise},
    \end{cases}
    \label{eq:pihm-G-entries}
\end{equation}
so that $f_Q'(x) = \braket{\tau(x)|\G|\psi}_n$, and the $m$-th derivative follows from the power $\G^{m}$. \Cref{eq:pihm-G-entries} is dense and upper triangular: row $k$ is non-zero at every column $p > k$ of the right parity, so the lowest-order coefficient can depend on as many as half of the others, and $\G$ carries $\Theta(N^2)$ non-zero entries in total. \Cref{app:cheb} gives this fill-in explicitly at $N=4$, alongside the banded factors $\Bt, \Dt$ that \Cref{sec:descriptor-construction} builds from the same recurrence.

\paragraph{The differential-equation row.} For a linear equation of order $k$ with constant coefficients, $\sum_{j=0}^k a_j f^{(j)} = 0$, direct substitution of $\G^{j}$ for each derivative gives
\begin{equation}
    \Astd = \sum_{j=0}^k a_j\, \G^{j},
    \label{eq:pihm-A}
\end{equation}
Wu et al.'s own construction \cite{wu2025pihm}; a second-order equation, for instance, gives $\Astd = a_2 \G^{2} + a_1 \G + a_0 I$ directly.

A coefficient $a_j(x)$ that varies with $x$ cannot be substituted this way: $a_j(x) f^{(j)}(x)$ is a product of two functions, and the Chebyshev coefficients of a product are not the product of the coefficients. Wu et al.\ resolve this with a second operator, the lifting operator $\M_{x^p}$, built from the general Chebyshev product-to-sum identity
\begin{equation}
    x^p\, T_k(x) = 2^{-p}\sum_{i=0}^{p}\binom{p}{i}\, T_{|k+p-2i|}(x),
    \label{eq:pihm-product-to-sum}
\end{equation}
the $p$-fold iterate of the pairwise rule $T_jT_k = \tfrac12(T_{j+k}+T_{|j-k|})$ that reappears below in the nonlinear extension. Conjugating \Cref{eq:pihm-product-to-sum} into the normalised bases of \Cref{eq:pihm-tau} gives an operator $\M_{x^p}$ satisfying
\begin{equation}
    x^p\,\bra{\tau(x)}_n = \bra{\tau(x)}_{n+1}\,\M_{x^p}
    \label{eq:pihm-Mxp}
\end{equation}
(entries: \Cref{app:pihm-residual}). Multiplying a degree-$(N-1)$ model by $x^p$ can raise its degree by as much as $p$, so $\M_{x^p}$ is exact and rectangular, $2N \times N$: no truncation is needed once the result is represented in the one-qubit-wider register \Cref{eq:pihm-Mxp} already targets. A general coefficient, expanded to a finite power series $a_j(x) = \sum_p c_{j,p}\,x^p$ (a Maclaurin truncation, for a non-polynomial coefficient), lifts by linearity to $\M_{a_j} = \sum_p c_{j,p}\, \M_{x^p}$, and the differential-equation row for the general, variable-coefficient case follows by composing $\M_{a_j}$ and $\G^{j}$ term-by-term,
\begin{equation}
    \Astd = \sum_{j=0}^k \M_{a_j}\, \G^{j},
    \label{eq:pihm-A-general}
\end{equation}
with $\M_{a_j} = a_j I$ recovering \Cref{eq:pihm-A} whenever $a_j$ is constant.
% REVISE: under the lift a constant coefficient enters as a_j M_1 (2N x N, the
%   paper's <tau|_{n+1} M_1), not a_j I (N x N); every term of one row must use
%   the same output register (Planning.md §3.3).

% LAND (new): sources. The inhomogeneous residual adds a collapse row; the
%   collapse row maps psi to the constant function f_q(x_s)/f(x_s), which is 1
%   on the true solution (Planning.md §3.2). The paper expands r(x) as a
%   Maclaurin series of order pbar = 3, 5, 7 (Figs 5b-d).
\begin{equation}
    \Astd = \sum_{j=0}^{k} \M_{a_j}\,\G^{j} - \M_r\,\Dzero(x_s),
    \qquad
    \Dzero(x_s) = \frac{\Brow(x_s)}{f(x_s)} .
    \label{eq:pihm-A-source}
\end{equation}

\paragraph{Data constraints.} Boundary and initial conditions are built from one rank-one point-evaluation generator, $\Brow_n(x) = \sqrt{2^n}\,\ket{0}^{\otimes n}\bra{\tau(x)}_n$, non-zero only in its first row. A condition with a zero target is \emph{invariant} -- unchanged under $\ket{\psi} \to c\ket{\psi}$ -- and enters the Hamiltonian directly as one further linear constraint, $\Brow\ket{\psi}=0$, e.g.\ a Dirichlet condition $f(x_z)=0$ giving $\Brow_n(x_z)\ket{\psi}=0$ and a Neumann condition $f'(x_m)=0$ giving $\Brow_n(x_m)\G\ket{\psi}=0$. A condition with a non-zero target, $f(x_s)=y_s$, is not scale-covariant and instead fixes the amplitude $\sqrt{\eta}$ of \Cref{eq:pihm-fq} directly, after the ground state has been found:
\begin{equation}
    \sqrt{\eta} = \frac{y_s}{\braket{\tau(x_s)|\psi_g}_n}.
    \label{eq:pihm-regular}
\end{equation}
Only one such \emph{regular} constraint is used per solve, and \Cref{eq:pihm-regular} is the only place $\eta$ re-enters the pipeline.
% REVISE: reconcile with eq:pihm-A-source -- the paper's regular point x_s also
%   supplies the collapse row D0(x_s) that injects sources (and lifts linear
%   terms in the doubled space), so a regular constraint does enter the residual
%   of an inhomogeneous equation.

\begin{definition}[Physics-informed effective Hamiltonian]
\label{def:pihm-hamiltonian}
Collecting the differential-equation row $\Astd$ and every invariant data constraint $\Brow_i$,
\begin{equation}
    \Hstd = \Astd^\top \Astd + \sum_i \Brow_i^\top \Brow_i, \qquad \braket{\psi|\Hstd|\psi} = \|\Astd\psi\|^2 + \sum_i \|\Brow_i\psi\|^2.
    \label{eq:pihm-hamiltonian}
\end{equation}
$\Hstd$ is Hermitian and positive semi-definite, so $\braket{\psi|\Hstd|\psi} \ge 0$ with equality exactly when every constraint holds simultaneously. The differential equation problem is thereby an eigenvalue problem, $\Hstd\ket{\psi_g} = \lambda_{\min}\ket{\psi_g}$, $\lambda_{\min}\approx0$, whose ground state $\ket{\psi_g}$ is the solution.
\end{definition}
% LAND: least squares, not an exact kernel -- lambda_min > 0 and small for a
%   non-polynomial solution; an exact kernel for a polynomial one (the paper's
%   Legendre panels). The ground state is the smallest right singular vector of
%   the stacked residual [A_std; B_1; ...] (Planning.md §3.3).

\subsection{The Published Circuits, Filter and Readout}
\label{sec:pihm-encoding}
% ~1 pp. Described AS PUBLISHED; no measured numbers of ours here (T-13).
% LAND, in order:
%   1. Figs 2 and 8: feature map + zero-state reflection S_{n+1} = 2|0><0| - I;
%      B on n + 2 qubits with alpha_B = sqrt(2^{n+1}); D0 with
%      alpha = sqrt(2^{n+1}) / f(x_s).
%   2. Fig. 9, the G ladder: an omega_c flag, an X ancilla, an n-qubit offset
%      register, a value ancilla, the system; U_R (n singly-controlled R_Y plus
%      n R-shifts) and U_C (n multi-controlled R_Y, one L-shift, 2n SWAPs); the
%      value load is a small-angle ladder, hence APPROXIMATE; the published
%      prefactor eq:pihm-published-prefactor. The paper's claims, stated
%      neutrally: 2n + 3 qubits, O(n^3) gates, entry error decreasing
%      exponentially in n (tested in §5.1 and §8.3).
%   3. Mixed-parity QSVT on H/||H||_F, and the paper's time rule
%      eq:pihm-time-rule (the paper writes t; tested in §8.4).
%   4. The interferometric readout (the paper's Eqs 36-39), described as
%      published (after kyriienko2020inverse, scali2023topological): the
%      all-zero probabilities of the prepared state and of an LCU-combined
%      state give the signed overlap tau(x).psi; eta from a regular datum.
% FIG (optional; drop if over budget): fig:pihm-published-pipeline.

\begin{equation}
    \alpha_{\G}^{\mathrm{pub}} = \bigl(2^{n-1} + 2^{n}\bigr)\,\|\G\|_S,
    \qquad
    \|\G\|_S = \max\bigl(\|\G\G^{\top}\|_1,\ \|\G^{\top}\G\|_1\bigr),
    \label{eq:pihm-published-prefactor}
\end{equation}

\begin{equation}
    \beta_{\mathrm{pub}} \ge \frac{\lambda_{\max}}{\lambda_2\,(n-1)} .
    \label{eq:pihm-time-rule}
\end{equation}

\placeholderfigure{The published pipeline of Wu et al.: residual, block encoding, QSVT imaginary-time filter and interferometric readout.}{fig:pihm-published-pipeline}{Optional schematic of the pipeline as published (Figs 2, 8 and 9 of the paper, QSVT-QITE on $H/\|H\|_F$, overlap readout), drawn to match the thesis's own pipeline figure (\Cref{fig:pipeline}) so the two can be compared at a glance.}

\subsection{The Doubled-Space Extension to Nonlinearity}
\label{sec:pihm-doubled}
% ~0.4 pp
% LAND: Psi = psi (x) psi with the fold (eq:pihm-fold); the paper's own
%   statement of >= 2^{2n-1} zero eigenvalues and that it "seeks a degenerate
%   state that closely matches the analytic solution". The argument that kernels
%   are subspaces and product states are not -- so the method PROJECTS rather
%   than prepares -- belongs here as ASSESSMENT; the formal statement is
%   prop:no-penalty (§7.2).

Wu et al.\ extend \Cref{def:pihm-hamiltonian} to polynomial nonlinearity by doubling the Hilbert space, $\ket{\Psi} = \ket{\psi}\otimes\ket{\psi}$, and folding a product of two coefficient vectors back into one register via the same product-to-sum identity, $T_jT_k = \tfrac12(T_{j+k}+T_{|j-k|})$ \cite{wu2025pihm}. A residual quadratic in $\ket{\psi}$ of this kind annihilates a subspace of the $2N$-dimensional $\ket{\Psi}$ far larger than the single physical direction: Wu et al.\ themselves note that the resulting $H$ carries \emph{at least} $2^{2n-1}$ degenerate zero-energy states, half of the doubled Hilbert space at minimum. No imaginary-time filter can select a direction inside a degenerate zero-energy cluster, so extracting the physical solution from this construction is classical post-selection within a known kernel rather than genuine state preparation; \Cref{sec:nonlinear} addresses this, for a restricted but practically important class of problems, with a different construction.
% REVISE: |Psi> = |psi> (x) |psi> lives in an N^2-dimensional space (2^{2n}),
%   not a 2N-dimensional one; "half of the doubled space" (2^{2n-1}) already
%   assumes N^2.

\begin{equation}
    \bra{\tau(x)}_n \otimes \bra{\tau(x)}_n = \bra{\tau(x)}_{n+1}\,\nfold_1 ,
    \qquad
    \nfold_1 \in \mathbb{R}^{2N \times N^2} .
    \label{eq:pihm-fold}
\end{equation}

\subsection{Critical Assessment}
\label{sec:pihm-assessment}
% ~0.6 pp -- THE PIVOT, and the HD-gate critical assessment.
% LAND, in order:
%   1. Strengths: a unified construction; exact kernels for polynomial
%      solutions; poly(n) gates; no discretisation grid.
%   2. Weaknesses, most severe first, each with a pointer that NAMES what is
%      measured and where (§8.2 sec:results-reproduction, §8.5
%      sec:results-scaling):
%        (1) the spectrum: ||H|| grows as N^8 against a flat gap, which fixes the
%            filter degree whatever the encoder;
%        (2) the published normalisation is loose by a factor growing like N^3;
%        (3) printed choices distort the solution at the printed n (rounded
%            constraint points, Maclaurin sources);
%        (4) the PDE invariant sets determine the solution only at the printed n;
%        (5) the nonlinear route cannot prepare;
%        (6) the readout identity assumes deterministic preparation.
%   3. Close with THREE ITALIC QUESTIONS, each answered by a named chapter:
%        (i)   Is the cost the encoder's or the Hamiltonian's?        -> sec:standard
%        (ii)  Can the residual be reformulated so that its Hamiltonian is banded
%              and well-conditioned, without changing its solution? -> sec:descriptor
%        (iii) Can nonlinear equations be prepared rather than projected?
%                                                                   -> sec:nonlinear
%   4. The spine sentence, as the question it answers (ThesisPlanning.md §3.2).

% CHECKLIST (marker)
%   1. Could any sentence be read as claiming authorship of the framework
%      (Style Trap 13)?
%   2. Is wu2025pihm cited at the head of sec:pihm-residual and at each
%      construction?
%   3. Are OUR numbers absent from §3.2-3.5, and present in §3.6 only as
%      pointed-to evidence?
%   4. Does §3.6 end with the three questions and the chapters that answer them?
%   5. Is the chapter <= 4 pp?
% TRAPS  "collocation"; the legacy adder-ladder reconstruction; the retired
%        Theta~(N^2) query floor.
